\documentclass[11pt,a4paper]{article}
\usepackage{fontspec}
\setmainfont{Noto Serif CJK SC}
\XeTeXlinebreaklocale "zh"
\XeTeXlinebreakskip = 0pt plus 1pt minus 0.1pt
\usepackage{geometry}
\geometry{left=2.3cm,right=2.3cm,top=2.2cm,bottom=2.2cm}
\usepackage{tabularx}
\usepackage{booktabs}
\usepackage{enumitem}
\setlength{\parindent}{0pt}
\setlength{\parskip}{0.55em}
\begin{document}
\begin{center}{\Large\bfseries 审稿意见回复}\end{center}
\vspace{0.2em}
\textbf{稿件：} 类型化判断与确定性执行：面向消息通道的同意门控智能体协作网关架构

\textbf{专题：} JCRAI 2026 专题 XII：智能体互联与协作智能

\textbf{作者：} 蔡亦张、刘军、马镝、禹可、李珂

\textbf{日期：} 2026 年 11 月 2 日

感谢审稿人认真而有建设性的阅读。这三条主要意见准确地指出了论文最薄弱的三个位置，我们已逐条在正文、图题与表格版式中处理。本轮为小修：只改动论文如何陈述其结果、以及让已有证据更容易核查，没有新增实验。

\section*{修订摘要}
\begin{center}\small
\begin{tabularx}{\textwidth}{@{}XX@{}}
\toprule
\textbf{修改} & \textbf{位置} \\
\midrule
对"回复文本"与"发现排名"的分歧做逐样例分析 & §9.1 \\
两幅图的图题补充消息流向与分支 & 图 1、图 2 \\
端口三元组在表题与表头中写明 & 表 5 \\
判断平面实验的设计写成 20 组对话 × 3 类事件 = 60 条 & §9.5 \\
删去与新增内容重复的句子，保持 15 页 & §9.6、§10 \\
对审稿意见提到的连词问题做了核对 & 全文 \\
\bottomrule
\end{tabularx}
\end{center}
\section*{主要问题一：回复文本抽取与 42/50 的一致性}
\begin{quote}\small §9.1 的"推荐"结果是从回复文本中解析出来的，而不是来自带埋点的 API（作者已承认）；它与发现排名的 42/50 一致性值得更深入的分析。\end{quote}
我们同意，回复解析是两种测量中较弱的一种。本轮利用已存记录做了逐样例分析，而不是只用一句话带过。

§9.1 现在报告三件事。其一，发现服务把目标排在第一位的有 40/50，进入前三的有 48/50；回复在三次投递中分别选中目标 48、48、50 条。其二，两种测量在 42/50 条上一致，且分歧是单向的——不存在"发现排第一、回复却漏掉"的情况。其三，第一次投递中那 8 条仅由回复命中的样本里，6 条在结构化列表中排第二或第三，2 条落在前三之外。因此，两种测量出现分歧时，回复是在同一候选池内重新排序或扩选，而不是换成了另一个候选池。

这些计数背后的逐样例记录是产物中的 \texttt{live\_\allowbreak{}discovery\_\allowbreak{}check\_\allowbreak{}v2.jsonl} 与 \texttt{live\_\allowbreak{}runs\_\allowbreak{}v2\_\allowbreak{}r1.jsonl} 至 \texttt{live\_\allowbreak{}runs\_\allowbreak{}v2\_\allowbreak{}r3.jsonl}，八条案例可逐条复核。§6 说明了如何从回复中读取选择，§9.1 说明了结构化排名对每条请求只读取一次。

\section*{主要问题二：只评测了一个托管模型}
\begin{quote}\small §9.3 只评测了一个托管模型（deepseek-flash）；类型化与自由文本的比较若要谈泛化，至少还需要一个模型。\end{quote}
这条属实，也是当前最主要的一项未完成工作。§9.3 开头就写明端点与设置——单一托管模型、温度 0、推理强度固定为低、不设输出上限——因此这项测量被限定为"在该端点下两种答案形态的表现"。摘要、结果与结论都没有作跨模型的主张。产物中包含套件、提示词与缓存回复，换一个端点即可在不改动测量代码的前提下重复同一比较。

本轮没有加入第二个模型。那需要对 280 条套件按同样三次重复重新运行，而本轮修订限定在"能用已有记录支撑的改动"之内。

\section*{主要问题三：模板化构造的注入压力集}
\begin{quote}\small §9.6 的注入压力集由模板生成，因此对真实攻击的鲁棒性尚未得到检验。\end{quote}
确实如此，论文本来就写明了构造方式。该集合是一个设计好的网格：8 个引导族 × 5 种载体 × 5 种权威框架 = 200 条注入项；8 个族 × 4 种载体 × 8 种报告形式 = 256 条良性近似项。盲复读对 456 条构造标签中的 64 条提出异议，仲裁改动了 51 条，最终留出集包含 67 条注入项与 47 条良性项。

这个压力集测量的是两种答案形态在"引用、转述、否认"这类文本上的误报行为，而不是对真实世界中收集到的攻击的鲁棒性。§9.3 描述了网格与仲裁过程，§9.6 现在用一句话写明模板构造，删去了周围的重复表述。

\section*{次要问题}
\textbf{图题。} 图 1 现在按消息流写明从入站（M1）、分区（M2）、发现（M3）、同意（M4）、执行（M5）、状态（M6）到投递（M7）的路径，并注明哪些判断点服务于哪个模块。图 2 写明了各分支：缓存命中直接返回、策略拒绝、决定待定时返回调用者并在用户答复后从认证一步重新进入、端点校验，以及失败分类。

\textbf{表 5 的端口写法。} 表题写明端口按"消融 / 同代码对照 / 线上实例"读取，表头使用同样的缩写，不必再从表题反推三元组的含义。

\textbf{§9.5 的样本量。} 设计现在集中写在一处：20 组对话、每组三类事件、共判定 60 条；0.9 阈值下复述召回 20/20（95\% Wilson [83.9\%, 100\%]），非复述 0/40 越阈（速率上界 7.5\%），判断调用中位 1.07 秒（0.53–1.74 秒）。我们保留了数字，把设计写清楚，而没有增加样本。

\textbf{连词与空格。} 我们核对了当前版本：第 1 页读作 "plane in advisory mode"，空格完整；对 PDF 文本做长词扫描也未发现粘连，两处命中是 \texttt{pdftotext} 换行造成的伪影，而不是 PDF 中的文字。如果审稿人手上的版本仍显示连写，值得检查一下所用的抽取工具。

\section*{建议}
\textbf{加入第二个判断模型。} 我们认可，这也是下一步最优先的一项。实验装置、套件与缓存提示词都在产物中，换一个端点即可重复同一比较。

\textbf{对推荐抽取做埋点。} 我们认可。本轮补充的是结构化测量本来就支持的排名层面分析；若要把回复解析换成协议层面的记录，需要改动适配器并重新投递 130 条请求，本轮未包含。

\textbf{报告判断平面的时延—准确率权衡。} 论文报告了 §9.5 的顾问调用时延与 §9.4 的网关侧处理开销，但没有跨阈值给出时延—准确率曲线。这条曲线需要在每个阈值上各跑一次，现有记录不包含这些数据。

\textbf{在图 2 中补上同意分支。} 已完成。图题现在描述了"决定待定"分支，以及用户答复后从认证一步重新进入。

\textbf{说明歧义请求的"组合"是否经过人工判定。} §9.1 已写明：组合只表示推荐了多于一个智能体，方案在语义上是否覆盖请求并未裁定。

\section*{对稿件的影响}
修订稿在 LNCS 版式下为 15 页，无 overfull box，无未定义引用。改动集中在 §9.1、图 1 与图 2 的图题、表 5 的表题与表头；§9.6 与结论删去了被这些新增内容重复的句子。

\end{document}
